% Options for packages loaded elsewhere
\PassOptionsToPackage{unicode}{hyperref}
\PassOptionsToPackage{hyphens}{url}
%
\documentclass[
  brazilian,
]{article}
\usepackage{lmodern}
\usepackage{amssymb,amsmath}
\usepackage{ifxetex,ifluatex}
\ifnum 0\ifxetex 1\fi\ifluatex 1\fi=0 % if pdftex
  \usepackage[T1]{fontenc}
  \usepackage[utf8]{inputenc}
  \usepackage{textcomp} % provide euro and other symbols
\else % if luatex or xetex
  \usepackage{unicode-math}
  \defaultfontfeatures{Scale=MatchLowercase}
  \defaultfontfeatures[\rmfamily]{Ligatures=TeX,Scale=1}
\fi
% Use upquote if available, for straight quotes in verbatim environments
\IfFileExists{upquote.sty}{\usepackage{upquote}}{}
\IfFileExists{microtype.sty}{% use microtype if available
  \usepackage[]{microtype}
  \UseMicrotypeSet[protrusion]{basicmath} % disable protrusion for tt fonts
}{}
\makeatletter
\@ifundefined{KOMAClassName}{% if non-KOMA class
  \IfFileExists{parskip.sty}{%
    \usepackage{parskip}
  }{% else
    \setlength{\parindent}{0pt}
    \setlength{\parskip}{6pt plus 2pt minus 1pt}}
}{% if KOMA class
  \KOMAoptions{parskip=half}}
\makeatother
\usepackage{xcolor}
\IfFileExists{xurl.sty}{\usepackage{xurl}}{} % add URL line breaks if available
\IfFileExists{bookmark.sty}{\usepackage{bookmark}}{\usepackage{hyperref}}
\hypersetup{
  pdftitle={Parte III: Processamento de Linguagem Natural para Linguistas},
  pdfauthor={CiberExt 26-29 · FEELT38103 · Universidade Federal de Uberlândia},
  pdflang={pt-BR},
  hidelinks,
  pdfcreator={LaTeX via pandoc}}
\urlstyle{same} % disable monospaced font for URLs
\usepackage{color}
\usepackage{fancyvrb}
\newcommand{\VerbBar}{|}
\newcommand{\VERB}{\Verb[commandchars=\\\{\}]}
\DefineVerbatimEnvironment{Highlighting}{Verbatim}{commandchars=\\\{\}}
% Add ',fontsize=\small' for more characters per line
\newenvironment{Shaded}{}{}
\newcommand{\AlertTok}[1]{\textcolor[rgb]{1.00,0.00,0.00}{\textbf{#1}}}
\newcommand{\AnnotationTok}[1]{\textcolor[rgb]{0.38,0.63,0.69}{\textbf{\textit{#1}}}}
\newcommand{\AttributeTok}[1]{\textcolor[rgb]{0.49,0.56,0.16}{#1}}
\newcommand{\BaseNTok}[1]{\textcolor[rgb]{0.25,0.63,0.44}{#1}}
\newcommand{\BuiltInTok}[1]{#1}
\newcommand{\CharTok}[1]{\textcolor[rgb]{0.25,0.44,0.63}{#1}}
\newcommand{\CommentTok}[1]{\textcolor[rgb]{0.38,0.63,0.69}{\textit{#1}}}
\newcommand{\CommentVarTok}[1]{\textcolor[rgb]{0.38,0.63,0.69}{\textbf{\textit{#1}}}}
\newcommand{\ConstantTok}[1]{\textcolor[rgb]{0.53,0.00,0.00}{#1}}
\newcommand{\ControlFlowTok}[1]{\textcolor[rgb]{0.00,0.44,0.13}{\textbf{#1}}}
\newcommand{\DataTypeTok}[1]{\textcolor[rgb]{0.56,0.13,0.00}{#1}}
\newcommand{\DecValTok}[1]{\textcolor[rgb]{0.25,0.63,0.44}{#1}}
\newcommand{\DocumentationTok}[1]{\textcolor[rgb]{0.73,0.13,0.13}{\textit{#1}}}
\newcommand{\ErrorTok}[1]{\textcolor[rgb]{1.00,0.00,0.00}{\textbf{#1}}}
\newcommand{\ExtensionTok}[1]{#1}
\newcommand{\FloatTok}[1]{\textcolor[rgb]{0.25,0.63,0.44}{#1}}
\newcommand{\FunctionTok}[1]{\textcolor[rgb]{0.02,0.16,0.49}{#1}}
\newcommand{\ImportTok}[1]{#1}
\newcommand{\InformationTok}[1]{\textcolor[rgb]{0.38,0.63,0.69}{\textbf{\textit{#1}}}}
\newcommand{\KeywordTok}[1]{\textcolor[rgb]{0.00,0.44,0.13}{\textbf{#1}}}
\newcommand{\NormalTok}[1]{#1}
\newcommand{\OperatorTok}[1]{\textcolor[rgb]{0.40,0.40,0.40}{#1}}
\newcommand{\OtherTok}[1]{\textcolor[rgb]{0.00,0.44,0.13}{#1}}
\newcommand{\PreprocessorTok}[1]{\textcolor[rgb]{0.74,0.48,0.00}{#1}}
\newcommand{\RegionMarkerTok}[1]{#1}
\newcommand{\SpecialCharTok}[1]{\textcolor[rgb]{0.25,0.44,0.63}{#1}}
\newcommand{\SpecialStringTok}[1]{\textcolor[rgb]{0.73,0.40,0.53}{#1}}
\newcommand{\StringTok}[1]{\textcolor[rgb]{0.25,0.44,0.63}{#1}}
\newcommand{\VariableTok}[1]{\textcolor[rgb]{0.10,0.09,0.49}{#1}}
\newcommand{\VerbatimStringTok}[1]{\textcolor[rgb]{0.25,0.44,0.63}{#1}}
\newcommand{\WarningTok}[1]{\textcolor[rgb]{0.38,0.63,0.69}{\textbf{\textit{#1}}}}
\setlength{\emergencystretch}{3em} % prevent overfull lines
\providecommand{\tightlist}{%
  \setlength{\itemsep}{0pt}\setlength{\parskip}{0pt}}
\setcounter{secnumdepth}{-\maxdimen} % remove section numbering
\ifxetex
  % Load polyglossia as late as possible: uses bidi with RTL langages (e.g. Hebrew, Arabic)
  \usepackage{polyglossia}
  \setmainlanguage[]{brazilian}
\else
  \usepackage[shorthands=off,main=brazilian]{babel}
\fi

\title{Parte III: Processamento de Linguagem Natural para Linguistas}
\usepackage{etoolbox}
\makeatletter
\providecommand{\subtitle}[1]{% add subtitle to \maketitle
  \apptocmd{\@title}{\par {\large #1 \par}}{}{}
}
\makeatother
\subtitle{Unidade 4 - Respostas dos Exercícios}
\author{CiberExt 26-29 · FEELT38103 · Universidade Federal de
Uberlândia}
\date{Agosto de 2026}

\begin{document}
\maketitle

\hypertarget{respostas-introduuxe7uxe3o-aos-word-embeddings}{%
\section{\texorpdfstring{Respostas --- Introdução aos \emph{word
embeddings}}{Respostas --- Introdução aos word embeddings}}\label{respostas-introduuxe7uxe3o-aos-word-embeddings}}

\hypertarget{parte-a-conceitos}{%
\subsection{Parte A --- Conceitos}\label{parte-a-conceitos}}

\textbf{Resposta 1.}

A hipótese distribucional afirma que \textbf{semelhança de significado
resulta em semelhança de distribuição linguística}: palavras que ocorrem
em contextos semelhantes tendem a ter significados semelhantes.

A citação de Firth a antecipa ao deslocar o critério de identificação do
significado da palavra isolada para sua \textbf{companhia}, isto é, para
os elementos com que ela coocorre. Firth não formula a hipótese em
termos distribucionais quantitativos --- isso é feito por Harris
(\href{https://doi.org/10.1080/00437956.1954.11659520}{1954}) ---, mas
estabelece o princípio metodológico de que o significado é contextual e,
portanto, observável no uso.

\textbf{Resposta 2.}

\begin{itemize}
\tightlist
\item
  O eixo \textbf{sintagmático} é o eixo da \textbf{combinação}: descreve
  como os elementos escolhidos se encadeiam horizontalmente para formar
  estruturas maiores (\emph{Helsinki} + \emph{is} + \emph{the capital
  of} + \emph{Finland}).
\item
  O eixo \textbf{paradigmático} é o eixo da \textbf{seleção}: descreve o
  conjunto de alternativas que poderiam ocupar uma mesma posição na
  estrutura (\emph{Helsinki} / \emph{Tallinn} / \emph{Estocolmo}). As
  alternativas se definem por oposição.
\end{itemize}

Na matriz da seção 2.1:

\begin{itemize}
\tightlist
\item
  Cada \textbf{coluna} corresponde a um lema único, isto é, a uma
  \textbf{escolha paradigmática} disponível.
\item
  Cada \textbf{linha} corresponde a um \emph{Doc}, isto é, a uma
  \textbf{estrutura sintagmática} concreta, que registra quais escolhas
  paradigmáticas foram efetivamente combinadas.
\end{itemize}

\textbf{Resposta 3.}

A representação é chamada de ``saco de palavras'' porque registra apenas
\textbf{quais} palavras ocorrem e \textbf{quantas vezes}, descartando
completamente a informação sobre a \textbf{ordem} em que aparecem. É
como se as palavras fossem jogadas em um saco e agitadas.

Duas limitações:

\begin{enumerate}
\def\labelenumi{\arabic{enumi}.}
\tightlist
\item
  \textbf{Perda da ordem e da estrutura.} As orações \emph{o cão mordeu
  o homem} e \emph{o homem mordeu o cão} recebem exatamente o mesmo
  vetor, embora signifiquem coisas opostas.
\item
  \textbf{Crescimento da dimensionalidade e esparsidade.} Cada palavra
  nova do vocabulário exige uma dimensão adicional. Como cada documento
  contém apenas uma fração ínfima do vocabulário, a esmagadora maioria
  das dimensões vale zero, o que desperdiça memória e dificulta o
  cálculo de distâncias significativas.
\end{enumerate}

\textbf{Resposta 4.}

Um vetor \textbf{esparso} é aquele em que a maior parte das dimensões
vale zero e apenas poucas carregam informação --- como o vetor one-hot,
em que exatamente uma dimensão vale 1. Um vetor \textbf{denso} é aquele
em que praticamente todas as dimensões têm valores não nulos, cada uma
codificando alguma parcela de informação.

Os embeddings da camada oculta são densos porque a rede é
\textbf{forçada} a comprimir a informação de 24 dimensões esparsas em
apenas 2 dimensões. Como não há espaço para reservar uma dimensão por
palavra, cada dimensão precisa passar a codificar informação distribuída
sobre a palavra e seu contexto.

\textbf{Resposta 5.}

O objetivo final não é dispor de um preditor de palavras vizinhas ---
essa predição, isoladamente, tem pouca utilidade prática. O objetivo é
obter \textbf{representações numéricas úteis} das palavras.

A predição funciona como tarefa substituta porque \textbf{só é possível
resolvê-la bem se a rede aprender boas representações}. Para prever que
\emph{the} e \emph{capital} aparecem perto de \emph{be}, a rede precisa
codificar, nos pesos da camada oculta, informação sobre o comportamento
distribucional de \emph{be}. Ao fim do treinamento, descartamos o
preditor e ficamos com os pesos --- que são os embeddings.

Trata-se do princípio geral do \textbf{aprendizado auto-supervisionado}:
o rótulo é extraído do próprio texto, sem anotação humana.

\textbf{Resposta 6.}

A camada oculta é um gargalo porque é muito mais estreita (2 neurônios)
do que as camadas de entrada e de saída (24 unidades cada). Toda a
informação necessária para a predição precisa \textbf{passar por essas
duas dimensões}.

Esse estrangulamento é justamente o que produz o aprendizado. Se a
camada oculta tivesse 24 neurônios, a rede poderia simplesmente copiar a
entrada, memorizando cada palavra individualmente, sem generalizar nada.
Ao ser obrigada a comprimir, a rede é forçada a \textbf{agrupar palavras
com comportamento distribucional semelhante} na mesma região do espaço
--- que é exatamente o que a hipótese distribucional prevê.

\hypertarget{parte-b-programauxe7uxe3o}{%
\subsection{Parte B --- Programação}\label{parte-b-programauxe7uxe3o}}

\textbf{Resposta 7.}

\begin{Shaded}
\begin{Highlighting}[]
\ImportTok{from}\NormalTok{ sklearn.metrics.pairwise }\ImportTok{import}\NormalTok{ cosine\_similarity}

\CommentTok{\# Compara os vetores dos Docs nos índices 3 e 4}
\NormalTok{cosine\_similarity([df.values[}\DecValTok{3}\NormalTok{]], [df.values[}\DecValTok{4}\NormalTok{]])}
\end{Highlighting}
\end{Shaded}

\begin{verbatim}
array([[0.37796447]])
\end{verbatim}

O valor é coerente. As orações são:

\begin{verbatim}
Travelling between Helsinki and Tallinn takes about two hours
Ferries depart from downtown Helsinki and Tallinn
\end{verbatim}

Elas compartilham três lemas --- \texttt{Helsinki}, \texttt{and} e
\texttt{Tallinn} ---, o que produz uma similaridade moderada (0,38).
Esse valor é claramente inferior ao dos \emph{Docs} 0 e 1 (0,67), que
compartilham quatro lemas em uma estrutura sintagmática idêntica
(\emph{is the capital of}), mas superior ao dos pares que quase nada
compartilham.

\textbf{Resposta 8.}

\begin{Shaded}
\begin{Highlighting}[]
\CommentTok{\# Acrescenta a nova oração à lista}
\NormalTok{examples.append(}\StringTok{"Stockholm is the capital of Sweden"}\NormalTok{)}

\CommentTok{\# Reprocessa toda a lista}
\NormalTok{docs }\OperatorTok{=} \BuiltInTok{list}\NormalTok{(nlp.pipe(examples))}

\CommentTok{\# Reconta os lemas}
\NormalTok{lemma\_counts }\OperatorTok{=}\NormalTok{ \{i: doc.count\_by(LEMMA) }\ControlFlowTok{for}\NormalTok{ i, doc }\KeywordTok{in} \BuiltInTok{enumerate}\NormalTok{(docs)\}}
\NormalTok{lemma\_counts }\OperatorTok{=}\NormalTok{ \{i: \{docs[i].vocab[k].text: v }\ControlFlowTok{for}\NormalTok{ k, v }\KeywordTok{in}\NormalTok{ counter.items()\}}
                \ControlFlowTok{for}\NormalTok{ i, counter }\KeywordTok{in}\NormalTok{ lemma\_counts.items()\}}

\CommentTok{\# Reconstrói o DataFrame}
\NormalTok{df }\OperatorTok{=}\NormalTok{ pd.DataFrame.from\_dict(lemma\_counts).sort\_index(ascending}\OperatorTok{=}\VariableTok{True}\NormalTok{).fillna(}\DecValTok{0}\NormalTok{).T}

\CommentTok{\# Verifica o formato}
\BuiltInTok{print}\NormalTok{(df.shape)}

\CommentTok{\# Recalcula a matriz de similaridade}
\NormalTok{sim }\OperatorTok{=}\NormalTok{ cosine\_similarity(df.values)}
\BuiltInTok{print}\NormalTok{(sim[}\DecValTok{5}\NormalTok{])}
\end{Highlighting}
\end{Shaded}

\begin{verbatim}
(6, 26)
[0.66666667 0.66666667 0.40824829 0.         0.         1.        ]
\end{verbatim}

O \emph{DataFrame} passa a ter \textbf{26 colunas}: as 24 originais mais
os dois lemas novos, \texttt{Stockholm} e \texttt{Sweden}.

Os \emph{Docs} 0 e 1 tornam-se \textbf{igualmente} os mais similares ao
novo \emph{Doc} 5, ambos com 0,67. Isso faz sentido: as três orações
compartilham exatamente a mesma estrutura sintagmática (\emph{X is the
capital of Y}), diferindo apenas nos dois nomes próprios. É precisamente
a previsão da hipótese distribucional: \texttt{Helsinki},
\texttt{Tallinn} e \texttt{Stockholm} são alternativas paradigmáticas de
uma mesma posição.

\textbf{Resposta 9.}

\begin{Shaded}
\begin{Highlighting}[]
\CommentTok{\# Similaridade entre \textquotesingle{}capital\textquotesingle{} e \textquotesingle{}the\textquotesingle{}}
\BuiltInTok{print}\NormalTok{(cosine\_similarity([lemma\_df[}\StringTok{\textquotesingle{}capital\textquotesingle{}}\NormalTok{].values], [lemma\_df[}\StringTok{\textquotesingle{}the\textquotesingle{}}\NormalTok{].values]))}

\CommentTok{\# Similaridade entre \textquotesingle{}Helsinki\textquotesingle{} e \textquotesingle{}ferry\textquotesingle{}}
\BuiltInTok{print}\NormalTok{(cosine\_similarity([lemma\_df[}\StringTok{\textquotesingle{}Helsinki\textquotesingle{}}\NormalTok{].values], [lemma\_df[}\StringTok{\textquotesingle{}ferry\textquotesingle{}}\NormalTok{].values]))}
\end{Highlighting}
\end{Shaded}

\begin{verbatim}
array([[0.53846154]])
array([[0.13363062]])
\end{verbatim}

\textbf{Comentário.} O par \texttt{capital} / \texttt{the} apresenta
similaridade alta (0,54) porque ambos ocorrem repetidamente nos mesmos
contextos locais (\emph{the capital of}), compartilhando vizinhos como
\texttt{be}, \texttt{of} e um ao outro. Isso ilustra um ponto
importante: a hipótese distribucional captura \textbf{similaridade de
distribuição}, que nem sempre corresponde a similaridade de significado
--- um determinante e um substantivo não são sinônimos.

O par \texttt{Helsinki} / \texttt{ferry} apresenta similaridade baixa
(0,13): as duas palavras ocorrem em vizinhanças bastante distintas no
nosso pequeno corpus.

\textbf{Resposta 10.}

\begin{Shaded}
\begin{Highlighting}[]
\CommentTok{\# Recria o DataFrame vazio}
\NormalTok{lemma\_df\_3 }\OperatorTok{=}\NormalTok{ pd.DataFrame(index}\OperatorTok{=}\NormalTok{unique\_lemmas, columns}\OperatorTok{=}\NormalTok{unique\_lemmas).fillna(}\DecValTok{0}\NormalTok{)}

\CommentTok{\# Percorre os Tokens, agora com janela de três palavras de cada lado}
\ControlFlowTok{for}\NormalTok{ token }\KeywordTok{in}\NormalTok{ combined\_docs:}

    \CommentTok{\# Note o intervalo ampliado de posições}
    \ControlFlowTok{for}\NormalTok{ position }\KeywordTok{in}\NormalTok{ [}\OperatorTok{{-}}\DecValTok{3}\NormalTok{, }\OperatorTok{{-}}\DecValTok{2}\NormalTok{, }\OperatorTok{{-}}\DecValTok{1}\NormalTok{, }\DecValTok{1}\NormalTok{, }\DecValTok{2}\NormalTok{, }\DecValTok{3}\NormalTok{]:}

        \ControlFlowTok{try}\NormalTok{:}
\NormalTok{            nbor\_lemma }\OperatorTok{=}\NormalTok{ token.nbor(position).lemma\_}
\NormalTok{            lemma\_df\_3.at[token.lemma\_, nbor\_lemma] }\OperatorTok{+=} \DecValTok{1}

        \ControlFlowTok{except} \PreprocessorTok{IndexError}\NormalTok{:}
            \ControlFlowTok{continue}

\CommentTok{\# Recalcula a similaridade}
\NormalTok{cosine\_similarity([lemma\_df\_3[}\StringTok{\textquotesingle{}Tallinn\textquotesingle{}}\NormalTok{].values], [lemma\_df\_3[}\StringTok{\textquotesingle{}Helsinki\textquotesingle{}}\NormalTok{].values])}
\end{Highlighting}
\end{Shaded}

\begin{verbatim}
array([[0.56577895]])
\end{verbatim}

\textbf{Comparação.} Com janela 2, a similaridade era de \textbf{0,43};
com janela 3, sobe para \textbf{0,57}. Ampliar a janela aumenta o número
de vizinhos compartilhados registrados e, portanto, o número de
dimensões não nulas em comum, elevando a similaridade.

De modo geral, janelas menores tendem a capturar relações mais
\textbf{sintáticas} (colocações imediatas), enquanto janelas maiores
capturam relações mais \textbf{temáticas} ou tópicas. A escolha da
janela é, portanto, uma decisão analítica, não um detalhe técnico.

\textbf{Resposta 11.}

\begin{Shaded}
\begin{Highlighting}[]
\ImportTok{import}\NormalTok{ numpy }\ImportTok{as}\NormalTok{ np}

\KeywordTok{def}\NormalTok{ one\_hot(lemma, vocabulary):}
    \CommentTok{"""Devolve o vetor one{-}hot de um lema em relação a um vocabulário.}

\CommentTok{    Parâmetros:}
\CommentTok{        lemma: string com o lema a ser codificado}
\CommentTok{        vocabulary: lista de lemas únicos}

\CommentTok{    Devolve:}
\CommentTok{        Um array NumPy do tamanho do vocabulário.}
\CommentTok{    """}

    \CommentTok{\# Verifica se o lema pertence ao vocabulário}
    \ControlFlowTok{if}\NormalTok{ lemma }\KeywordTok{not} \KeywordTok{in}\NormalTok{ vocabulary:}

        \CommentTok{\# Levanta um erro informativo caso contrário}
        \ControlFlowTok{raise} \PreprocessorTok{ValueError}\NormalTok{(}\SpecialStringTok{f"O lema \textquotesingle{}}\SpecialCharTok{\{}\NormalTok{lemma}\SpecialCharTok{\}}\SpecialStringTok{\textquotesingle{} não está no vocabulário."}\NormalTok{)}

    \CommentTok{\# Cria um vetor de zeros com o tamanho do vocabulário}
\NormalTok{    vector }\OperatorTok{=}\NormalTok{ np.zeros(shape}\OperatorTok{=}\BuiltInTok{len}\NormalTok{(vocabulary))}

    \CommentTok{\# Define como 1 a posição correspondente ao lema}
\NormalTok{    vector[vocabulary.index(lemma)] }\OperatorTok{=} \DecValTok{1}

    \ControlFlowTok{return}\NormalTok{ vector}


\CommentTok{\# Testes}
\BuiltInTok{print}\NormalTok{(one\_hot(}\StringTok{\textquotesingle{}Helsinki\textquotesingle{}}\NormalTok{, unique\_lemmas))}
\BuiltInTok{print}\NormalTok{(one\_hot(}\StringTok{\textquotesingle{}Curitiba\textquotesingle{}}\NormalTok{, unique\_lemmas))}
\end{Highlighting}
\end{Shaded}

\begin{verbatim}
[0. 1. 0. 0. 0. 0. 0. 0. 0. 0. 0. 0. 0. 0. 0. 0. 0. 0. 0. 0. 0. 0. 0. 0.]
ValueError: O lema 'Curitiba' não está no vocabulário.
\end{verbatim}

\textbf{Resposta 12.}

\begin{Shaded}
\begin{Highlighting}[]
\CommentTok{\# Reconstrói a rede com quatro neurônios na camada oculta}
\NormalTok{network\_input }\OperatorTok{=}\NormalTok{ keras.Input(shape}\OperatorTok{=}\NormalTok{targets.shape[}\DecValTok{1}\NormalTok{], name}\OperatorTok{=}\StringTok{\textquotesingle{}input\_layer\textquotesingle{}}\NormalTok{)}

\CommentTok{\# Apenas o argumento \textquotesingle{}units\textquotesingle{} muda}
\NormalTok{hidden\_layer }\OperatorTok{=}\NormalTok{ Dense(units}\OperatorTok{=}\DecValTok{4}\NormalTok{, activation}\OperatorTok{=}\VariableTok{None}\NormalTok{, name}\OperatorTok{=}\StringTok{\textquotesingle{}hidden\_layer\textquotesingle{}}\NormalTok{)(network\_input)}

\NormalTok{output\_layer }\OperatorTok{=}\NormalTok{ Dense(units}\OperatorTok{=}\NormalTok{context.shape[}\DecValTok{1}\NormalTok{], activation}\OperatorTok{=}\StringTok{\textquotesingle{}softmax\textquotesingle{}}\NormalTok{,}
\NormalTok{                     name}\OperatorTok{=}\StringTok{\textquotesingle{}output\_layer\textquotesingle{}}\NormalTok{)(hidden\_layer)}

\NormalTok{model\_4d }\OperatorTok{=}\NormalTok{ keras.Model(inputs}\OperatorTok{=}\NormalTok{network\_input, outputs}\OperatorTok{=}\NormalTok{output\_layer)}
\NormalTok{model\_4d.}\BuiltInTok{compile}\NormalTok{(loss}\OperatorTok{=}\StringTok{\textquotesingle{}categorical\_crossentropy\textquotesingle{}}\NormalTok{)}

\NormalTok{model\_4d.fit(x}\OperatorTok{=}\NormalTok{targets, y}\OperatorTok{=}\NormalTok{context, batch\_size}\OperatorTok{=}\DecValTok{64}\NormalTok{, epochs}\OperatorTok{=}\DecValTok{1500}\NormalTok{, verbose}\OperatorTok{=}\DecValTok{0}\NormalTok{)}

\CommentTok{\# Verifica o formato dos pesos da camada oculta}
\NormalTok{weights\_4d }\OperatorTok{=}\NormalTok{ model\_4d.get\_weights()[}\DecValTok{0}\NormalTok{]}
\BuiltInTok{print}\NormalTok{(weights\_4d.shape)}
\end{Highlighting}
\end{Shaded}

\begin{verbatim}
(24, 4)
\end{verbatim}

A matriz de pesos passa a ter formato \texttt{(24,\ 4)}: cada um dos 24
lemas é agora representado por um vetor de \textbf{quatro} dimensões.

Já não é possível plotar o espaço diretamente porque um gráfico de
dispersão só dispõe de dois eixos, e teríamos quatro coordenadas por
ponto. Para visualizar, seria necessário aplicar uma técnica de
\textbf{redução de dimensionalidade} --- como PCA ou t-SNE --- que
projeta os vetores de quatro dimensões em duas, preservando
aproximadamente as distâncias relativas. Essa é exatamente a abordagem
adotada na próxima unidade, ao visualizar os embeddings do spaCy, que
têm centenas de dimensões.

\textbf{Resposta 13.}

\begin{Shaded}
\begin{Highlighting}[]
\CommentTok{\# Localiza os índices dos dois lemas}
\NormalTok{i\_hel }\OperatorTok{=}\NormalTok{ unique\_lemmas.index(}\StringTok{\textquotesingle{}Helsinki\textquotesingle{}}\NormalTok{)}
\NormalTok{i\_tal }\OperatorTok{=}\NormalTok{ unique\_lemmas.index(}\StringTok{\textquotesingle{}Tallinn\textquotesingle{}}\NormalTok{)}

\CommentTok{\# Compara os embeddings aprendidos}
\NormalTok{cosine\_similarity([hidden\_layer\_weights[i\_hel]], [hidden\_layer\_weights[i\_tal]])}
\end{Highlighting}
\end{Shaded}

\begin{verbatim}
array([[0.60469307]])
\end{verbatim}

\textbf{Comparação.} A similaridade calculada a partir da matriz de
coocorrência era de \textbf{0,43}; a partir dos embeddings aprendidos,
obtemos aproximadamente \textbf{0,60}.

Duas ressalvas importantes:

\begin{enumerate}
\def\labelenumi{\arabic{enumi}.}
\tightlist
\item
  O valor exato \textbf{varia entre execuções}, pois os pesos são
  inicializados aleatoriamente. O que se deve observar é a tendência,
  não o número.
\item
  Em apenas duas dimensões, quaisquer dois vetores tendem a apresentar
  similaridade de cossenos relativamente alta, simplesmente porque há
  pouco espaço para que fiquem distantes. Comparar valores absolutos
  entre espaços de dimensionalidades diferentes é, portanto, enganoso.
\end{enumerate}

\textbf{Resposta 14.}

\begin{enumerate}
\def\labelenumi{\arabic{enumi}.}
\item
  \textbf{Vocabulário insuficiente para generalizar.} Com 24 lemas e 118
  pares, quase todas as palavras ocorrem uma ou duas vezes. A rede não
  tem evidência estatística suficiente para distinguir regularidade de
  acaso: uma única coocorrência tem o mesmo peso que um padrão
  sistemático.
\item
  \textbf{Sobreajuste ao exemplo específico.} O modelo aprende as cinco
  orações particulares, não a língua inglesa. \texttt{Helsinki} e
  \texttt{Tallinn} ficam próximos porque aparecem literalmente lado a
  lado nos dados, e não porque o modelo tenha inferido que ambos são
  capitais bálticas.
\item
  \textbf{Ausência de palavras funcionais em contexto variado.} Palavras
  como \texttt{the} e \texttt{of} ocorrem em pouquíssimos contextos
  distintos, quando na realidade sua característica definidora é ocorrer
  em uma enorme variedade deles. Suas representações ficam, portanto,
  distorcidas --- e, como são frequentes, distorcem também as
  representações das palavras vizinhas.
\end{enumerate}

Uma quarta consequência, mais sutil: as próprias métricas de avaliação
tornam-se pouco informativas, já que não há dados retidos para teste.

\textbf{Resposta 15.}

O modelo é incapaz de fazer a distinção porque atribui
\textbf{exatamente um vetor a cada forma lexical}. O dicionário
\texttt{lemma\_vectors} mapeia a \emph{string}
\texttt{\textquotesingle{}banco\textquotesingle{}} a um único vetor
one-hot, e a camada oculta produz um único embedding correspondente. As
duas acepções de ``banco'' são, para o modelo, literalmente a mesma
entrada --- e o vetor resultante acaba sendo uma média confusa dos dois
contextos de uso.

Esse é o limite dos chamados embeddings \textbf{estáticos} (Word2vec,
GloVe, fastText): um tipo lexical, um vetor.

Para superá-lo é necessário produzir embeddings \textbf{contextuais}, em
que a representação de uma palavra depende da sentença em que ela
ocorre. Isso exige:

\begin{itemize}
\tightlist
\item
  Uma arquitetura capaz de \textbf{ler a sentença inteira} ao codificar
  cada palavra --- tipicamente um \emph{Transformer} com mecanismo de
  atenção, em vez de uma rede que recebe uma única palavra isolada como
  entrada.
\item
  Uma tarefa de treinamento que \textbf{force o uso do contexto}, como
  prever uma palavra mascarada a partir de toda a sentença (o objetivo
  de \emph{masked language modelling} usado pelo BERT).
\item
  Volume muito maior de dados e de parâmetros, para que os diferentes
  contextos de uso de cada forma sejam efetivamente observados.
\end{itemize}

Nesses modelos, as ocorrências de ``banco'' em \emph{sentei no banco} e
em \emph{fui ao banco} recebem vetores distintos, pois cada
representação é computada em função das palavras que a cercam. Os
embeddings contextuais obtidos de \emph{Transformers} são o tema da
próxima unidade.

\end{document}
